\chapter{DistributionFitter}
\label{chap:distribution_fitter}

\section{Purpose}

The \texttt{DistributionFitter} class is used to fit one or more candidate
probability distributions to an observed dataset and to compare them using
goodness-of-fit metrics. In actuarial practice, this step typically
precedes simulation: the analyst inspects historical claim counts or
severities, selects a plausible distributional form, and validates the
choice using diagnostic statistics.

\section{Supported Distributions}

\texttt{DistributionFitter} supports the same distributions exposed by
\texttt{actstats}, summarised below.

\begin{center}
\begin{tabular}{ll}
\toprule
\textbf{Type} & \textbf{Distributions} \\
\midrule
Continuous severity & \texttt{normal}, \texttt{logistic},
                      \texttt{exponential}, \texttt{gamma}, \texttt{beta},
                      \texttt{lognormal}, \texttt{weibull},
                      \texttt{pareto}, \texttt{uniform} \\
Discrete frequency  & \texttt{poisson}, \texttt{negative binomial} \\
\bottomrule
\end{tabular}
\end{center}

The distributions used in any particular fitting run can be controlled
through the \texttt{distributions} argument or via a YAML configuration file.

\section{Goodness-of-Fit Metrics}

For each fitted distribution, \texttt{DistributionFitter} computes:

\begin{itemize}[leftmargin=*]
    \item Log-likelihood
    \item Akaike Information Criterion (AIC)
    \item Bayesian Information Criterion (BIC)
    \item Chi-square statistic
    \item Kolmogorov--Smirnov statistic
\end{itemize}

The information criteria are defined as
\[
    \mathrm{AIC} = 2k - 2\ln L, \qquad
    \mathrm{BIC} = k \ln n - 2 \ln L,
\]
where $k$ is the number of estimated parameters, $n$ is the sample size,
and $L$ is the likelihood evaluated at the maximum-likelihood estimate.

By default, the metrics used for ranking are AIC and BIC. Additional
metrics can be requested by passing a custom \texttt{metrics} list.

\section{Basic Workflow}

The typical workflow uses four steps:

\begin{enumerate}[leftmargin=*]
    \item Load configuration and select candidate distributions.
    \item Construct a \texttt{DistributionFitter} with the data.
    \item Call \texttt{fit()} to estimate parameters and compute metrics.
    \item Review and select the best-fit distribution.
\end{enumerate}

\begin{lstlisting}[style=pythonstyle]
from actsim import DistributionFitter, load_config
from actstats import actuarial as act

# Generate example severity data
sev_data = act.lognormal(0.5, 0.2).rvs(size=10000)

# Load configuration
config = load_config()
distributions = config.distributions["severity"]
metrics = config.metrics

# Construct the fitter
fitter = DistributionFitter(
    sev_data,
    distributions=distributions,
    metrics=metrics,
)

# Fit candidate distributions
fitter.fit()
\end{lstlisting}

\section{Inspecting Results}

After fitting, the best-fit distribution under each metric is stored in
the \texttt{best\_fits} attribute, while the currently selected
distribution is stored in \texttt{selected\_fit}.

\begin{lstlisting}[style=pythonstyle]
print("Best fits:", fitter.best_fits)
print("Selected fit:", fitter.selected_fit)
print("Selected distribution:", fitter.get_selected_dist())
print("Selected parameters:", fitter.get_selected_params())
\end{lstlisting}

\section{Manually Selecting a Distribution}

Users can override the automatic selection by name:

\begin{lstlisting}[style=pythonstyle]
fitter.select_distribution("gamma")
print(fitter.selected_fit["params"])
\end{lstlisting}

This is useful when the analyst has a prior reason to prefer a particular
distribution, or when comparing scenarios.

\section{Sampling and Prediction}

Once a distribution is selected, the fitter can be used as a drop-in
sampler. The \texttt{sample} method returns random draws from the
selected distribution, while \texttt{sample\_mixed} inserts a proportion
of zeros and ones to handle datasets with mass at those points (e.g.\
zero-inflated severity data).

\begin{lstlisting}[style=pythonstyle]
# Draw 10 samples
samples = fitter.sample(size=10)

# Draw 100 samples with 10% zeros and 5% ones
mixed = fitter.sample_mixed(zero_prop=0.1, one_prop=0.05, size=100)

# Evaluate the predicted density at given points
density = fitter.predict([0.5, 1.0, 1.5, 2.0])
\end{lstlisting}

\section{Diagnostic Plots}

The \texttt{plot\_predictions()} method draws a histogram of the data
together with the probability density functions of the fitted candidate
distributions. This plot is the primary visual diagnostic for assessing
fit quality.

\begin{lstlisting}[style=pythonstyle]
fitter.plot_predictions()
\end{lstlisting}

\section{Summary Reports}

Two methods produce numerical summaries:

\begin{itemize}[leftmargin=*]
    \item \texttt{calculate\_statistics()} returns a comparison of the
          mean, standard deviation, and key percentiles between the
          observed data and the predicted density.
    \item \texttt{summary()} returns a DataFrame of all fitted
          distributions, their parameters, and their goodness-of-fit
          metrics.
\end{itemize}

\begin{lstlisting}[style=pythonstyle]
fitter.calculate_statistics()
fitter.summary()
\end{lstlisting}

\section{Practical Notes}

A few practical considerations apply when using \texttt{DistributionFitter}:

\begin{itemize}[leftmargin=*]
    \item Some distributions (e.g.\ Pareto, beta) require positive or
          bounded support. If the input data violate the support, the
          fit for that distribution may fail; \texttt{DistributionFitter}
          catches such failures and continues with the remaining
          distributions.
    \item AIC and BIC are convenient ranking tools but should be
          interpreted alongside diagnostic plots, especially in the
          tails, which often matter most for actuarial applications.
    \item For frequency data, choose distributions appropriate to count
          data (e.g.\ Poisson, negative binomial) rather than continuous
          distributions.
\end{itemize}
